\chapter{Oligopoly: Cournot and Bertrand}\label{ch:b1:cournot-and-bertrand}

An isolated town has one dentist. When the town grows, a second dentist
opens, then a third. How does the price of a filling change as each one
arrives, and how much larger must the town be to support a fourth? A
monopoly and a competitive industry are the two ends; between them are
markets with a few sellers, each of whom knows that its choices move the
price the others get. This chapter studies two models of such markets ---
sellers choosing quantities and sellers choosing prices --- and what the
number of sellers does in each.

\begin{openingchart}{Town population needed per firm to support $N$ firms,
five professions in isolated US towns, with the Cournot benchmark for
identical firms, scaled so that one firm needs 1 thousand people.}{Bresnahan and Reiss 1991, Table 5;
Cournot benchmark computed}
\begin{tikzpicture}
\begin{axis}[omrecord, xlabel={Number of firms $N$}, ylabel={Thousand people per firm}, xmin=0.8, xmax=5.6, ymin=0, ymax=2.2,
  xtick={1,...,5}, width=0.7\linewidth, height=0.42\linewidth, legend style={font=\footnotesize, at={(0.5,-0.24)},
  anchor=north, legend columns=3}]
  \addplot+[thick] table[x=n, y=doctors, col sep=comma] {figdata/bachelor-1/17-cournot-and-bertrand/out/per_firm.csv};
  \addplot+[thick] table[x=n, y=dentists, col sep=comma] {figdata/bachelor-1/17-cournot-and-bertrand/out/per_firm.csv};
  \addplot+[thick] table[x=n, y=druggists, col sep=comma] {figdata/bachelor-1/17-cournot-and-bertrand/out/per_firm.csv};
  \addplot+[thick] table[x=n, y=plumbers, col sep=comma] {figdata/bachelor-1/17-cournot-and-bertrand/out/per_firm.csv};
  \addplot+[thick] table[x=n, y=tire_dealers, col sep=comma] {figdata/bachelor-1/17-cournot-and-bertrand/out/per_firm.csv};
  \addplot[black, dashed, thick, mark=none] table[x=n, y=cournot, col sep=comma] {figdata/bachelor-1/17-cournot-and-bertrand/out/per_firm.csv};
  \legend{Doctors, Dentists, Druggists, Plumbers, Tire dealers, Cournot}
\end{axis}
\end{tikzpicture}
\end{openingchart}

A monopoly dentist needed a town of $0.71$ thousand people; two dentists
needed $2.54$ thousand, $1.27$ thousand per dentist. % ledger: b1:bresnahan1991-dentists
If each new firm only split a fixed pie of profit, the population per firm
would stay flat. It rises instead, and steeply from one firm to two: the
per-firm threshold almost doubles for doctors ($1.98$ times), dentists
($1.79$), druggists ($2.00$) and tire dealers ($1.82$). % figdata: ratios
A second firm must cut margins, so each firm needs more customers to cover
its fixed costs. After the third firm the lines are almost flat: the fourth
and fifth firms change little. The dashed line shows what identical
firms competing in quantities would give: from one firm to two its
per-firm threshold rises by a factor of only 1.125, far less in proportion
than the data's.

\begin{recall}
Oligopoly, the Nash \omterm{def:b1:models-and-data:equilibrium}{equilibrium} (a profile of strategies in which each
player's strategy is a best reply to the others') and the prisoner's
dilemma were met at school. Monopoly and the \omterm{def:b1:monopoly-pricing:lerner}{Lerner index} are in
\cref{ch:b1:monopoly-pricing}.
\end{recall}

\section{Competition in quantities}\label{sec:b1:cournot-and-bertrand:cournot}

Two firms make the same good at constant marginal cost $c$ and choose
quantities $q_1$, $q_2$ at the same time; the price is then the one at which
buyers take $Q = q_1 + q_2$, $p = a - bQ$.

\begin{definition}[Best-response function]\label{def:b1:cournot-and-bertrand:best-response}
\index{best-response function}
A firm's \emph{best-response function} gives, for each choice of its rivals,
the choice that maximises its own profit.
\end{definition}

Firm 1's profit is $(a - b(q_1 + q_2) - c)q_1$; setting its derivative to zero gives
\[
q_1 = \frac{a - c - bq_2}{2b},
\]
so firm 1 produces less when firm 2 produces more.

\begin{definition}[Cournot equilibrium]\label{def:b1:cournot-and-bertrand:cournot}
\index{Cournot equilibrium}
A \emph{Cournot \omterm{def:b1:models-and-data:equilibrium}{equilibrium}} is a list of quantities, one per firm, each a
best response to the others': a Nash \omterm{def:b1:models-and-data:equilibrium}{equilibrium} of the game in which
firms choose quantities.
\end{definition}

\begin{proposition}[Cournot with $n$ firms]\label{prop:b1:cournot-and-bertrand:n-firms}
With $n$ identical firms, linear demand $p = a - bQ$ and marginal cost $c$, the
\omterm{def:b1:cournot-and-bertrand:cournot}{Cournot equilibrium} has $q_i = \dfrac{a - c}{b(n + 1)}$ and
$p = c + \dfrac{a - c}{n + 1}$. Each firm's \omterm{def:b1:monopoly-pricing:lerner}{Lerner index} is $L = -1/(n\varepsilon)$, with $\varepsilon$
the elasticity of \omterm{def:b1:consumer-choice:market-demand}{market demand}.
\end{proposition}

\begin{proof}
Firm $i$'s condition is $a - bQ - bq_i - c = 0$. In a symmetric \omterm{def:b1:models-and-data:equilibrium}{equilibrium}
$Q = nq_i$, so $a - c = b(n + 1)q_i$, and $p = a - bnq_i = c + (a - c)/(n + 1)$. More
generally, firm $i$'s condition is $p + q_i\,\dd p/\dd Q = c$, so
$(p - c)/p = -\frac{q_i}{Q}\frac Qp\frac{\dd p}{\dd Q} = -s_i/\varepsilon$, with $s_i = 1/n$ its share.
\end{proof}

\begin{proposition}[The competitive limit]\label{prop:b1:cournot-and-bertrand:limit}
As $n \to \infty$, the Cournot price tends to marginal cost, and the margin
$p - c = (a - c)/(n + 1)$ falls like $1/n$.
\end{proposition}

With one firm the margin is half of $a - c$; with two, a third; with nine, a
tenth. The model gives the same result as monopoly at $n = 1$ and as
competition in the limit.

\begin{omfigure}
\begin{tikzpicture}
\begin{axis}[omrecord, grid=none, xlabel={$q_1$}, ylabel={$q_2$}, xmin=0, xmax=95, ymin=0, ymax=95,
  width=0.46\linewidth, height=0.44\linewidth]
  \addplot[demand, domain=0:90] {90 - 2*x};
  \addplot[supply, domain=0:90] {(90 - x)/2};
  \addplot[only marks, mark=*, mark size=1.8pt] coordinates {(30,30)};
  \addplot[only marks, mark=o, mark size=1.8pt] coordinates {(45,0) (0,45) (90,0) (0,90)};
  \addplot[guide] coordinates {(45,0) (0,45)};
  \addplot[guide] coordinates {(90,0) (0,90)};
  \node[font=\footnotesize, anchor=south west] at (axis cs:30,30) {Cournot};
  \node[font=\footnotesize, anchor=west] at (axis cs:4,72) {firm 1's best response};
  \node[font=\footnotesize, anchor=south west] at (axis cs:74,11) {firm 2's};
  \node[font=\footnotesize, anchor=south west] at (axis cs:46,1) {monopoly};
  \node[font=\footnotesize, anchor=south west] at (axis cs:64,27) {competition};
\end{axis}
\end{tikzpicture}

{\small Demand $p = 120 - Q$ and $c = 30$. The best responses
$q_1 = (90 - q_2)/2$ (blue) and $q_2 = (90 - q_1)/2$ (red) cross at 30 each. The dotted lines
are the outputs that sum to the monopoly total, 45, and the competitive
total, 90.}
\end{omfigure}

\section{Competition in prices}\label{sec:b1:cournot-and-bertrand:bertrand}

Now let the two firms choose prices. Buyers buy from the cheaper firm,
and split equally at equal prices; each firm can serve the whole market.

\begin{definition}[Bertrand equilibrium]\label{def:b1:cournot-and-bertrand:bertrand}
\index{Bertrand equilibrium}
A \emph{Bertrand \omterm{def:b1:models-and-data:equilibrium}{equilibrium}} is a list of prices, one per firm, each a best
response to the others'.
\end{definition}

\begin{proposition}[Two are enough]\label{prop:b1:cournot-and-bertrand:bertrand}
With two firms with the same marginal cost $c$, a homogeneous good and no
limit on capacity, the only \omterm{def:b1:cournot-and-bertrand:bertrand}{Bertrand equilibrium} is $p_1 = p_2 = c$.
\end{proposition}

\begin{proof}
At $p_1 = p_2 = c$ profits are zero; raising a price sells nothing, cutting it
sells below cost: no firm gains. Any other profile has a profitable move.
If the lower price exceeds $c$, the firm with the higher price (or either
firm, if they are equal) can take the whole market by charging slightly
less than the lower price, gaining almost the whole market's profit instead
of half or nothing. If one price equals $c$ and the other is higher, the
firm at $c$ earns nothing and gains by raising its price to just below the
other's. If a price is below $c$, the firm charging it loses money on every
unit and can raise it.
\end{proof}

\begin{remark}
\index{Bertrand paradox}
The \emph{Bertrand paradox} is the gap between this prediction and what is
observed: two sellers are enough for marginal-cost pricing in the model,
while in most real markets two sellers earn margins well above cost.
\end{remark}

With unequal costs $c_1 < c_2$, the low-cost firm prices just below $c_2$ and
takes the market, earning $c_2 - c_1$ a unit: only a cost advantage brings a
margin.

\section{Capacities and differentiation}\label{sec:b1:cournot-and-bertrand:soften}

Two things undo the paradox. The first is capacity. If firms first choose
capacities and then set prices, a firm cannot take the whole market by
undercutting, because it cannot serve it; under conditions studied in
Year~2, the outcome is the Cournot one. Quantity competition is then a model
of firms that commit to capacity --- seats on a ferry, rooms in a hotel ---
and compete in prices afterwards.

The second is differentiation. If the goods differ --- in location, in
brand, in quality --- a firm that charges a little more keeps some of its
buyers. With demands $q_i = \alpha - \beta p_i + \gamma p_j$, firm $i$'s best response to
$p_j$ is $p_i = (\alpha + \beta c + \gamma p_j)/(2\beta)$: it raises its price when its rival does.

\begin{definition}[Strategic substitutes and complements]\label{def:b1:cournot-and-bertrand:strategic}
\index{strategic substitutes}\index{strategic complements}
Choices are \emph{strategic substitutes} when a firm's best response falls
as its rival's choice rises, as with Cournot quantities, and
\emph{strategic complements} when it rises, as with differentiated Bertrand
prices.
\end{definition}

\begin{omfigure}
\begin{tikzpicture}
\begin{axis}[omrecord, grid=none, xlabel={$p_1$}, ylabel={$p_2$}, xmin=0, xmax=70, ymin=0, ymax=70,
  width=0.46\linewidth, height=0.44\linewidth]
  \addplot[demand, domain=30:47.5] {4*x - 120};
  \addplot[supply, domain=0:70] {(120 + x)/4};
  \addplot[only marks, mark=*, mark size=1.8pt] coordinates {(40,40)};
  \node[font=\footnotesize, anchor=north west] at (axis cs:40,40) {equilibrium};
  \node[font=\footnotesize, anchor=east] at (axis cs:44,62) {firm 1};
  \node[font=\footnotesize, anchor=south] at (axis cs:60,46) {firm 2};
\end{axis}
\end{tikzpicture}

{\small Differentiated price competition with $q_i = 100 - 2p_i + p_j$ and a
marginal cost of 10: the best responses $p_i = (120 + p_j)/4$ slope up and cross at
$p_1 = p_2 = 40$, well above cost.}
\end{omfigure}

\section{The leader and the follower}\label{sec:b1:cournot-and-bertrand:stackelberg}

Suppose firm 1 chooses its quantity first, and firm 2 sees it before
choosing. Firm 2 then plays its best response, and firm 1, knowing this,
chooses its quantity with the response built in: the game is solved
backwards.

\begin{definition}[Stackelberg equilibrium]\label{def:b1:cournot-and-bertrand:stackelberg}
\index{Stackelberg equilibrium}\index{leader--follower model}
In a \emph{Stackelberg} (leader--follower) \emph{\omterm{def:b1:models-and-data:equilibrium}{equilibrium}} the leader
chooses its quantity to maximise its profit given the follower's
\omterm{def:b1:cournot-and-bertrand:best-response}{best-response function}, and the follower then plays its best response.
\end{definition}

\begin{proposition}[The leader's advantage]\label{prop:b1:cournot-and-bertrand:leader}
With linear demand $p = a - bQ$ and a common marginal cost $c$, the leader
produces $(a - c)/(2b)$, the monopoly quantity, the follower half as much, and
the leader earns more than in the \omterm{def:b1:cournot-and-bertrand:cournot}{Cournot equilibrium}.
\end{proposition}

\begin{proof}
Given $q_2 = (a - c - bq_1)/(2b)$, the leader's profit is
$(a - b q_1 - b q_2 - c)q_1 = \tfrac12(a - c - bq_1)q_1$, which is largest at $q_1 = (a - c)/(2b)$; then
$q_2 = (a - c)/(4b)$. The leader's profit is $(a - c)^2/(8b)$, against $(a - c)^2/(9b)$ in
Cournot.
\end{proof}

The leader gains by committing: once its output is fixed, the follower
must accommodate it. A capacity built before the rival arrives is such a
commitment.

\section{Margins and the number of competitors}\label{sec:b1:cournot-and-bertrand:entry}

\begin{definition}[Entry threshold]\label{def:b1:cournot-and-bertrand:entry-threshold}
\index{entry threshold}
The \emph{entry threshold} $S_N$ is the smallest market size --- the number of
buyers --- at which $N$ firms can each cover their fixed costs.
\end{definition}

With identical Cournot firms, demand $Q = S(a - p)$ for a market of $S$ people,
marginal cost $c$ and fixed cost $F$, each of $N$ firms earns
$S(a - c)^2/(N + 1)^2 - F$, so $S_N = F(N + 1)^2/(a - c)^2$. The per-firm threshold $S_N/N$ rises
because margins fall as firms enter; if margins did not fall, it would be
flat. Comparing towns of different sizes thus measures how fast
competition erodes margins without observing a single price or cost.

\begin{record}[Entry in isolated towns]
Using towns in the United States far from any larger market, a study
counted the doctors, dentists, druggists, plumbers and tire dealers in
each and estimated the population needed for each number of firms. A
monopoly tire dealer or druggist needs about 500 people, a doctor or
dentist 700 to 900, a plumber about 1,400 (twice a dentist). % ledger: b1:bresnahan1991, b1:bresnahan1991-plumbers
In markets with five or fewer firms almost all of the change in
competitive conduct comes with the second and third firms (Bresnahan and
Reiss 1991). % ledger: b1:bresnahan1991
\end{record}

\begin{record}[Two are few]
Laboratory markets let the number of sellers be chosen. In Cournot
markets with two to five sellers, two sellers colluded somewhat, three
produced at the Cournot level, and four or five were never collusive, some
close to the competitive outcome (Huck, Normann and Oechssler 2004). % ledger: b1:huck2004
In Bertrand markets the prediction of marginal-cost pricing did poorly
with two sellers but, after some learning, well with three or four
(Dufwenberg and Gneezy 2000). % ledger: b1:dufwenberg2000
\end{record}

\begin{model}{Cournot oligopoly}
\assumes $n$ identical firms choosing quantities (or capacities) at the
same time; a homogeneous good.
\predicts Margins above cost falling like $1/n$; a \omterm{def:b1:monopoly-pricing:lerner}{Lerner index} of $-s_i/\varepsilon$;
\omterm{def:b1:cournot-and-bertrand:entry-threshold}{entry thresholds} per firm rising with $n$.
\tested Per-firm \omterm{def:b1:cournot-and-bertrand:entry-threshold}{entry thresholds} rise with the number of firms in
isolated towns, faster from one to two firms than Cournot with linear
demand predicts (Bresnahan and Reiss 1991); % ledger: b1:bresnahan1991
laboratory markets with three firms produce at the Cournot level (Huck,
Normann and Oechssler 2004). % ledger: b1:huck2004
\breaks When firms collude (Year~2); when they compete in prices with
spare capacity.
\end{model}

\begin{model}{Bertrand competition}
\assumes Firms setting prices, a homogeneous good, no limits on capacity,
equal costs.
\predicts Price equal to marginal cost with two firms or more, and no
effect of further entry.
\tested In the laboratory, prices near cost appear with three or four
sellers but not with two (Dufwenberg and Gneezy 2000). % ledger: b1:dufwenberg2000
\breaks In most real markets two sellers earn margins: products differ,
capacities bind, and the game is repeated.
\end{model}

\section*{Exercises}\label{sec:b1:cournot-and-bertrand:exercises}

\begin{exercise}[$\star$]\label{exo:b1:cournot-and-bertrand:duo}
Find the \omterm{def:b1:cournot-and-bertrand:cournot}{Cournot equilibrium} of two firms with $p = 120 - Q$ and $c = 30$: outputs,
price and profits.
\end{exercise}

\begin{exercise}[$\star$]\label{exo:b1:cournot-and-bertrand:bert}
Two firms with marginal cost 30 compete in prices for the same buyers.
What is the \omterm{def:b1:models-and-data:equilibrium}{equilibrium} price and profit?
\end{exercise}

\begin{exercise}[$\star$]\label{exo:b1:cournot-and-bertrand:read}
From Bresnahan and Reiss's Table 5, the source of the opening chart, a
monopoly dentist needs a town of $0.71$
thousand people and two dentists $2.54$ thousand. % ledger: b1:bresnahan1991-dentists
Compute the per-firm thresholds and their ratio.
\end{exercise}

\begin{exercise}[$\star$]\label{exo:b1:cournot-and-bertrand:third}
In the market of \cref{exo:b1:cournot-and-bertrand:duo}, predict the price
when a third identical firm enters.
\end{exercise}

\begin{exercise}[$\star\star$]\label{exo:b1:cournot-and-bertrand:surplus}
In the same market with $n$ Cournot firms, write the price and \omterm{def:b1:partial-equilibrium-welfare:surplus}{total
surplus} as functions of $n$, and evaluate for $n = 1, 2, 3$ and the limit.
\end{exercise}

\begin{exercise}[$\star\star$]\label{exo:b1:cournot-and-bertrand:share}
Show that a Cournot firm with market share $s$, possibly different from its
rivals', has \omterm{def:b1:monopoly-pricing:lerner}{Lerner index} $s/|\varepsilon|$.
\end{exercise}

\begin{exercise}[$\star\star$]\label{exo:b1:cournot-and-bertrand:unequal}
Two firms with marginal costs 10 and 12 compete in prices for a market with
demand $Q = 100 - p$. Find the \omterm{def:b1:models-and-data:equilibrium}{equilibrium} (prices in whole cents, or as a limit).
\end{exercise}

\begin{exercise}[$\star\star$]\label{exo:b1:cournot-and-bertrand:stack}
Solve the \omterm{def:b1:cournot-and-bertrand:stackelberg}{leader--follower model} for the market of
\cref{exo:b1:cournot-and-bertrand:duo}, and compare outputs and profits
with the Cournot outcome.
\end{exercise}

\begin{exercise}[$\star\star$]\label{exo:b1:cournot-and-bertrand:diff}
Firms with marginal cost 10 face $q_i = 100 - 2p_i + p_j$. Find the \omterm{def:b1:cournot-and-bertrand:bertrand}{Bertrand
equilibrium} prices, quantities and profits.
\end{exercise}

\begin{exercise}[$\star\star\star$]\label{exo:b1:cournot-and-bertrand:limit-ex}
Prove that the Cournot price tends to marginal cost as $n \to \infty$, and that
$n(p - c)$ tends to a limit; find it.
\end{exercise}

\begin{exercise}[$\star\star\star$]\label{exo:b1:cournot-and-bertrand:thresholds}
With demand $Q = S(a - p)$, cost $c$ and fixed cost $F$, derive $S_N$ for Cournot
firms and the ratios $S_{N+1}/S_N$ and $(S_{N+1}/(N + 1))/(S_N/N)$. What would
Bertrand competition imply for $S_2$?
\end{exercise}

\begin{exercise}[$\star\star\star$]\label{exo:b1:cournot-and-bertrand:refute}
What evidence would refute ``the second firm brings the price to marginal
cost''? Say what the \omterm{def:b1:cournot-and-bertrand:entry-threshold}{entry thresholds} of the opening chart would look like
if it were true.
\end{exercise}

\section{Problem: Ferries to the Islands}\label{sec:b1:cournot-and-bertrand:problem}

\begin{problem}[{Weekend problem --- how big an island three ferries need}]\label{pb:b1:cournot-and-bertrand:ferries}
Brenland's islands (invented) are served by ferries. An island of $S$
thousand people demands $Q = S(60 - p)$ thousand trips a year at a fare of $p$
dalens. A ferry's marginal cost is 20 a trip and its fixed cost 800 thousand
dalens a year. Money below is in thousands of dalens a year.

\medskip\noindent\textbf{Part I --- One ferry.} Island A has 4 thousand people.
\begin{enumerate}
  \item Write the inverse demand.
  \item Find the monopoly fare and trips.
  \item Compute the ferry's profit after its fixed cost.
  \item Compute the \omterm{def:b1:monopoly-pricing:lerner}{Lerner index} and the elasticity of demand at that fare.
  \item Compute \omterm{def:b1:consumer-welfare:consumer-surplus}{consumer surplus}.
  \item Compute the \omterm{def:b1:partial-equilibrium-welfare:dwl}{deadweight loss}.
\end{enumerate}

\medskip\noindent\textbf{Part II --- Two ferries choose capacities.} Island B has 9
thousand people and two ferries, each choosing how many trips to offer.
\begin{enumerate}[resume]
  \item Write the inverse demand.
  \item Find ferry 1's best response to ferry 2's trips.
  \item Find the \omterm{def:b1:cournot-and-bertrand:cournot}{Cournot equilibrium}: trips per ferry, in all, and the fare.
  \item Compute each ferry's profit.
  \item Compute the \omterm{def:b1:monopoly-pricing:lerner}{Lerner index} and check it against $-1/(n\varepsilon)$.
  \item Compute \omterm{def:b1:consumer-welfare:consumer-surplus}{consumer surplus}.
  \item Compare with a single ferry on island B: fare and profit. Would the two ferries gain by merging?
\end{enumerate}

\medskip\noindent\textbf{Part III --- A price war.}
\begin{enumerate}[resume]
  \item If the two ferries had ample capacity and set fares, what fare would result?
  \item What would each ferry's profit be, and what would happen next?
  \item If one ferry left, what would the other earn?
  \item At the Cournot fare each ferry runs full. Show that neither gains by cutting its fare, and that ferry 1 does not gain by offering fewer trips.
  \item Ferry 1 chooses its capacity before ferry 2 arrives. Find both capacities, the fare and the profits.
  \item What is ferry 1's commitment worth?
\end{enumerate}

\medskip\noindent\textbf{Part IV --- How many ferries?}
\begin{enumerate}[resume]
  \item Write the profit of each of $n$ Cournot ferries on an island of $S$ thousand people.
  \item Find the \omterm{def:b1:cournot-and-bertrand:entry-threshold}{entry thresholds} $S_1$, $S_2$ and $S_3$.
  \item Compute the per-ferry thresholds and their ratios.
  \item How many ferries can island A and island B support, and what is island B's fare then?
  \item Compare the first per-ferry ratio with Bresnahan and Reiss's ratios for doctors and dentists, and say what it suggests.
  \item What would $S_2$ be under Bertrand competition with ample capacity?
  \item State the result: the island size that supports three ferries, as a multiple of the size that supports one.
\end{enumerate}
\end{problem}
